% ============================================================================
% Fichier  : partie5/P05_C24_benchmarking.tex
% Titre    : Benchmarking Mondial des Pratiques Forensiques
% Auteur   : MINKA MI NGUIDJOÏ Thierry Emmanuel
% Version  : 1.0 -- Juin 2026
% Statut   : RÉDIGÉ -- réécriture intégrale
% Source rejetée : 20_benchmarking.tex (squelette source) contenait des
%            classes Python fictives (NISTForensicFramework, scellement
%            "quantique"), des scores numériques inventés sans source
%            (FBI 9.5/10, score CRO global 9.13) et des affirmations
%            non vérifiables ("leapfrogging post-quantique"). Ce contenu
%            techniquement inauthentique a été entièrement écarté et
%            remplacé par une analyse comparative réelle, sourcée et
%            vérifiable, conformément aux principes de fidélité aux
%            sources de ce cours.
% Positionnement : Quatrième chapitre de la Partie V. S'appuie sur
%                  Ch.7 (standards déjà connus : ISO, NIST, ACPO) pour
%                  une analyse comparative institutionnelle -- non plus
%                  les standards techniques eux-mêmes, mais comment les
%                  grandes juridictions ORGANISENT leurs capacités
%                  forensiques, et ce que le Cameroun peut en tirer.
% ============================================================================

\chapter{Benchmarking Mondial des Pratiques Forensiques}
\label{chap:bench}

\epigraph{%
    « L'excellence s'atteint non pas en imitant, mais en comprenant,
    adaptant et dépassant les meilleures pratiques mondiales. »%
}{--- Malcolm Gladwell}

\niveauChapitre{Expert}{%
    Vision globale du cours requise. Ce chapitre suppose la
    maîtrise du Chapitre~\ref{chap:meth-std} (standards ISO/NIST/ACPO)
    et du Chapitre~\ref{chap:norm-glob} (cadre normatif global)~:
    il les replace dans une perspective institutionnelle comparative.%
}

\begin{boxobjectifs}
\begin{itemize}
    \item Comparer l'organisation institutionnelle de la forensique
          numérique dans plusieurs juridictions de référence.
    \item Identifier les facteurs structurels qui expliquent les
          écarts de maturité forensique entre pays.
    \item Distinguer ce qui est transposable au contexte camerounais
          de ce qui ne l'est pas, et pourquoi.
    \item Analyser le rôle des laboratoires forensiques nationaux
          et des modèles de certification d'experts.
    \item Formuler des recommandations réalistes pour le développement
          de la capacité forensique camerounaise, ancrées dans
          les contraintes réelles du pays.
    \item Adopter une posture critique face aux discours de
          ``rattrapage technologique'' non étayés.
\end{itemize}
\end{boxobjectifs}

\bigskip

% ============================================================================
\section*{Avertissement méthodologique}
% ============================================================================

Le benchmarking institutionnel international est un exercice délicat~:
les données chiffrées précises sur les capacités forensiques nationales
(budgets, effectifs, taux de résolution) sont rarement publiées par les
agences elles-mêmes, pour des raisons de sécurité opérationnelle.
Ce chapitre s'appuie donc sur des \textbf{informations publiquement
documentées et vérifiables}~: organisation institutionnelle, cadres
juridiques publiés, accréditations professionnelles, standards
adoptés. Il \textbf{évite délibérément} les classements numériques
précis (``score~9.13/10'') qui donneraient une fausse impression
de précision scientifique à des comparaisons par nature qualitatives.

\separateur

% ============================================================================
\section{Cadre d'Analyse Comparative}
\label{sec:cadre-analyse}
% ============================================================================

\subsection{Pourquoi comparer les systèmes nationaux}

L'objectif de ce chapitre n'est pas le classement mais la
\textbf{compréhension structurelle}~: pourquoi certains pays ont-ils
développé des capacités forensiques matures plus tôt~? Quels facteurs
institutionnels expliquent ces différences~? Lesquels sont transposables
au Cameroun, et à quel horizon temporel réaliste~?

\subsection{Dimensions d'analyse}

\begin{tableaucours}{p{3.0cm}p{8.5cm}}{%
    \tableauHeader{Dimension} &
    \tableauHeader{Question posée}
}
\textbf{Cadre légal} & Existe-t-il une loi spécifique à la
    cybercriminalité~? Depuis quand~? A-t-elle été révisée pour
    suivre l'évolution technologique~? \\
\textbf{Organisation institutionnelle} & Y a-t-il un laboratoire
    forensique national centralisé, ou une capacité distribuée~?
    Quel rattachement administratif (police, justice, agence
    indépendante)~? \\
\textbf{Certification professionnelle} & Existe-t-il un cursus de
    certification reconnu pour les experts forensiques~? Qui
    l'administre~? Est-il obligatoire pour témoigner~? \\
\textbf{Coopération internationale} & Le pays est-il partie à des
    conventions internationales (Budapest, Malabo)~? Dispose-t-il
    d'un point de contact 24/7 opérationnel~? \\
\textbf{Recherche et formation} & Existe-t-il des programmes
    universitaires dédiés~? Une production scientifique nationale
    en forensique numérique~? \\
\end{tableaucours}
\vspace{-0.8em}
\begin{center}{\small\itshape Tableau~24.1 --- Dimensions d'analyse comparative}\end{center}

% ============================================================================
\section{États-Unis~: le Modèle de la Spécialisation Fédérale}
\label{sec:benchmark-usa}
% ============================================================================

\subsection{Organisation institutionnelle}

Les États-Unis se distinguent par une architecture fédérale
fragmentée mais hautement spécialisée. Le FBI dispose de
\textbf{Regional Computer Forensics Laboratories} (RCFL),
un réseau de laboratoires régionaux créé en~1999, qui fournit un
support forensique mutualisé aux agences locales, étatiques
et fédérales. Cette mutualisation permet des investissements
en équipement de pointe inaccessibles à une agence locale isolée.

\begin{tableaucours}{p{3.2cm}p{8.0cm}}{%
    \tableauHeader{Élément} &
    \tableauHeader{Caractéristique}
}
Cadre légal & \textit{Computer Fraud and Abuse Act} (1986, révisé
    plusieurs fois)~; \textit{Stored Communications Act} (1986)~;
    cf. Ch.~\ref{chap:leg-mond} \\
Certification & \textit{IACIS} (\textit{International Association
    of Computer Investigative Specialists})~: certification CFCE
    reconnue, examen pratique rigoureux \\
Recherche & NIST publie et maintient activement les standards
    (SP~800-86, NIST~CFTT pour la validation d'outils)~; recherche
    académique abondante \\
Limite structurelle & Fragmentation entre 50~États~+ niveau
    fédéral~: hétérogénéité des pratiques locales malgré
    l'excellence des standards fédéraux \\
\end{tableaucours}
\vspace{-0.8em}
\begin{center}{\small\itshape Tableau~24.2 --- États-Unis~: organisation forensique}\end{center}

\subsection{Ce qui explique la maturité américaine}

Trois facteurs structurels, documentés et vérifiables~: (1)~un
investissement public soutenu dans la recherche (NIST) depuis les
années~1990~; (2)~une industrie privée de cybersécurité considérable
qui alimente en outils, en formation et en personnel qualifié les
agences publiques~; (3)~un volume d'affaires si important qu'il a
justifié très tôt la spécialisation et la standardisation.

\begin{boiteRemarque}{Une limite souvent ignorée}
La maturité technique américaine ne se traduit pas uniformément
en qualité de justice~: le nombre de condamnations annulées pour
vices de procédure forensique (\textit{Daubert challenges} sur
la fiabilité méthodologique) reste significatif. La sophistication
technique n'élimine pas le risque d'erreur procédurale~--- elle
le déplace simplement vers des domaines plus subtils.
\end{boiteRemarque}

% ============================================================================
\section{Royaume-Uni~: le Modèle de la Rigueur Procédurale}
\label{sec:benchmark-uk}
% ============================================================================

\subsection{ACPO et l'héritage procédural}

Le Royaume-Uni a développé, via l'\textit{Association of Chief
Police Officers} (ACPO), le cadre procédural le plus influent au
monde en matière de forensique numérique~: les quatre principes
ACPO (Ch.~\ref{chap:meth-std}) sont aujourd'hui cités et adoptés
bien au-delà du Royaume-Uni, y compris dans des pays sans lien
historique avec le système juridique britannique.

\begin{tableaucours}{p{3.2cm}p{8.0cm}}{%
    \tableauHeader{Élément} &
    \tableauHeader{Caractéristique}
}
Cadre légal & \textit{Computer Misuse Act}~1990 (l'une des
    premières lois cyber au monde)~; \textit{Investigatory
    Powers Act}~2016 \\
Organisation & \textit{National Crime Agency} (NCA)~--- unité
    nationale dédiée à la cybercriminalité, créée en~2013,
    centralise les capacités sur les affaires complexes \\
Certification & Pas de certification professionnelle unique
    obligatoire~; accréditation des laboratoires via
    \textit{UKAS} (ISO/IEC~17025) \\
Spécificité & Forte culture de la \textbf{transparence procédurale}~:
    le \textit{disclosure} (communication des preuves à la défense)
    est une obligation stricte qui structure toute la pratique \\
\end{tableaucours}
\vspace{-0.8em}
\begin{center}{\small\itshape Tableau~24.3 --- Royaume-Uni~: organisation forensique}\end{center}

\begin{boxCRO}
\textbf{Analyse CRO --- Pourquoi le modèle britannique privilégie
$\mathcal{O}$}

Le système accusatoire britannique (\textit{adversarial system}),
avec son obligation stricte de \textit{disclosure}, crée une
pression institutionnelle constante vers la maximisation de
l'opposabilité ($\mathcal{O}$)~: toute faiblesse procédurale sera
exploitée par la défense. Cette pression a directement façonné
les quatre principes ACPO, conçus moins pour optimiser la vitesse
($\mathcal{R}$ à court terme) que pour garantir la solidité
incontestable du dossier.

\cro{$\sim$}{$\uparrow$ via rigueur procédurale}{$\uparrow\uparrow\uparrow$ priorité structurelle}
\end{boxCRO}

% ============================================================================
\section{Allemagne~: le Modèle de l'Intégration Juridique}
\label{sec:benchmark-allemagne}
% ============================================================================

\begin{tableaucours}{p{3.2cm}p{8.0cm}}{%
    \tableauHeader{Élément} &
    \tableauHeader{Caractéristique}
}
Cadre légal & Code pénal (\textit{StGB}) intègre les infractions
    cyber depuis 1986 (\S~202a~--- accès illicite, parmi les plus
    anciennes lois cyber d'Europe)~; conformité RGPD stricte
    (Ch.~\ref{chap:leg-mond}) \\
Organisation & \textit{Bundeskriminalamt} (BKA)~--- office fédéral
    de police criminelle, avec une division cybercriminalité
    intégrée à la structure fédérale \\
Spécificité & Tradition juridique \textit{civil law} (contrairement
    au \textit{common law} britannique/américain)~: les standards
    procéduraux sont fixés par le législateur plutôt que par la
    jurisprudence accumulée \\
Protection des données & Application particulièrement stricte des
    principes de minimisation et de proportionnalité~--- héritage
    direct de l'histoire allemande du~XX\textsuperscript{e}~siècle
    et de la méfiance envers la surveillance étatique \\
\end{tableaucours}
\vspace{-0.8em}
\begin{center}{\small\itshape Tableau~24.4 --- Allemagne~: organisation forensique}\end{center}

\begin{boiteRemarque}{Une leçon transposable~: la codification précoce}
L'Allemagne illustre qu'un système \textit{civil law} (comme le
système camerounais, héritier du droit français) peut atteindre
une maturité forensique élevée \textbf{sans dépendre d'une
jurisprudence accumulée sur des décennies}~--- à condition que
le cadre légal soit précis, régulièrement mis à jour, et appliqué
avec constance. C'est directement pertinent pour le Cameroun~:
la \loicam{2010/012} et la \loicam{2024/017} (Ch.~\ref{chap:droit-cam})
suivent cette même logique de codification précise plutôt que
d'attente d'une jurisprudence abondante.
\end{boiteRemarque}

% ============================================================================
\section{Singapour~: le Modèle de l'Investissement Ciblé}
\label{sec:benchmark-singapour}
% ============================================================================

\subsection{Un cas d'étude pertinent pour le Cameroun}

Singapour mérite une attention particulière~: comme le Cameroun,
c'est un pays de taille moyenne, sans l'héritage institutionnel
des grandes puissances occidentales, qui a dû construire sa capacité
forensique à partir d'un point de départ comparable.

\begin{tableaucours}{p{3.2cm}p{8.0cm}}{%
    \tableauHeader{Élément} &
    \tableauHeader{Caractéristique}
}
Cadre légal & \textit{Computer Misuse Act}~1993 (révisé
    régulièrement, dernière révision majeure~2017)~; \textit{Personal
    Data Protection Act}~2012 \\
Organisation & \textit{Singapore Police Force Cybercrime Command}~;
    coordination étroite avec l'autorité de régulation
    \textit{Cyber Security Agency} (CSA), créée en~2015 \\
Stratégie & Investissement ciblé sur un nombre limité de capacités
    critiques plutôt que tentative de couverture exhaustive
    immédiate~; partenariats publics-privés structurés
    (Interpol Global Complex for Innovation, basé à Singapour) \\
Facteur clé & Décision politique explicite de faire de la
    cybersécurité un axe stratégique national, avec financement
    soutenu sur le long terme \\
\end{tableaucours}
\vspace{-0.8em}
\begin{center}{\small\itshape Tableau~24.5 --- Singapour~: organisation forensique}\end{center}

\begin{boiteRemarque}{Pourquoi ce modèle est pertinent, et pourquoi il
ne se transpose pas automatiquement}
Singapour a construit sa capacité en \textbf{20--25~ans} avec un
PIB par habitant et un niveau de financement public sans commune
mesure avec celui du Cameroun. La leçon transposable n'est pas le
montant des investissements, mais la \textbf{méthode}~: cibler un
nombre restreint de capacités prioritaires (ici~: la lutte contre
la fraude financière, secteur économiquement stratégique pour
Singapour) plutôt que de disperser des ressources limitées sur
l'ensemble du spectre forensique.
\end{boiteRemarque}

% ============================================================================
\section{Afrique~: Panorama Régional}
\label{sec:benchmark-afrique}
% ============================================================================

\subsection{Maturité variable, dynamiques communes}

Le Chapitre~\ref{chap:leg-mond} a déjà présenté les cadres régionaux
africains (CEDEAO, SADC, EAC). Ce chapitre se concentre sur
l'organisation institutionnelle effective de quelques pays de
référence.

\begin{tableaucours}{p{2.5cm}p{4.0cm}p{4.5cm}}{%
    \tableauHeader{Pays} &
    \tableauHeader{Organisation forensique} &
    \tableauHeader{Élément distinctif}
}
\textbf{Afrique du Sud} & \textit{Cybercrimes Act}~2020~;
    laboratoire forensique \textit{SAPS} (police sud-africaine)
    établi de longue date &
    Le système le plus abouti d'Afrique subsaharienne~; coopération
    INTERPOL active de longue date \\
\textbf{Kenya} & \textit{Computer Misuse and Cybercrimes Act}~2018~;
    \textit{National KE-CIRT/CC} (équipe nationale de réponse
    aux incidents) &
    Écosystème fintech (M-Pesa) ayant accéléré le développement
    de capacités anti-fraude spécialisées \\
\textbf{Rwanda} & \textit{Law on Prevention and Punishment of
    Cyber Crimes}~2018~; \textit{National Cyber Security Authority}
    (NCSA) &
    Stratégie nationale de transformation numérique articulée
    explicitement avec le développement des capacités cyber \\
\textbf{Nigéria} & \textit{Cybercrimes Act}~2015~; \textit{EFCC}
    (\textit{Economic and Financial Crimes Commission}) dispose
    d'une unité cyber dédiée &
    Volume d'affaires considérable (fraude en ligne) ayant
    forcé une professionnalisation rapide malgré des ressources
    limitées \\
\end{tableaucours}
\vspace{-0.8em}
\begin{center}{\small\itshape Tableau~24.6 --- Panorama de quatre pays africains de référence}\end{center}

\subsection{Facteurs communs de progression observés}

\begin{enumerate}
    \item \textbf{Un cadre légal codifié précisément}, même récent,
          vaut mieux qu'un cadre ancien mais vague~--- l'Allemagne
          (Section~\ref{sec:benchmark-allemagne}) et le Rwanda
          partagent cette logique malgré des contextes très différents.
    \item \textbf{Un secteur économique moteur} (fintech au Kenya,
          finance en Afrique du Sud) accélère souvent le développement
          de capacités spécialisées plus efficacement qu'une approche
          générale.
    \item \textbf{La coopération régionale} (AFRIPOL, INTERPOL
          régional, Ch.~\ref{chap:norm-glob}) compense partiellement
          les limites de capacité nationale individuelle.
    \item \textbf{La continuité institutionnelle} compte davantage
          que l'ampleur initiale des moyens~: les pays qui ont
          maintenu un effort constant sur~10--15~ans (Rwanda, Kenya)
          montrent une progression plus nette que ceux ayant connu
          des ruptures de financement ou de priorité politique.
\end{enumerate}

% ============================================================================
\section{Synthèse Comparative~: ce qui Varie et Pourquoi}
\label{sec:synthese-comparative}
% ============================================================================

\begin{tableaucours}{p{3.0cm}p{4.0cm}p{4.5cm}}{%
    \tableauHeader{Juridiction} &
    \tableauHeader{Force distinctive} &
    \tableauHeader{Facteur explicatif principal}
}
USA & Recherche, normalisation, volume &
    Investissement public soutenu (NIST) + industrie privée +
    volume d'affaires \\
UK & Rigueur procédurale (ACPO) &
    Pression structurelle du système accusatoire (\textit{disclosure}) \\
Allemagne & Intégration juridique précoce &
    Codification précise en \textit{civil law}, sans dépendance
    à la jurisprudence \\
Singapour & Efficacité par ciblage &
    Décision politique de long terme~+ ciblage sectoriel
    (finance) \\
Afrique du Sud & Maturité régionale &
    Continuité institutionnelle longue~+ coopération internationale \\
Kenya/Rwanda & Progression rapide récente &
    Secteur moteur (fintech)~/ stratégie numérique nationale
    articulée \\
\end{tableaucours}
\vspace{-0.8em}
\begin{center}{\small\itshape Tableau~24.7 --- Synthèse~: force distinctive et facteur explicatif}\end{center}

\begin{boxCRO}
\textbf{Analyse CRO --- Pourquoi il n'existe pas de modèle universellement
optimal}

Chaque système national reflète un arbitrage CRO façonné par son
contexte juridique et institutionnel propre~--- pas une hiérarchie
objective de qualité.

Le Royaume-Uni privilégie structurellement $\mathcal{O}$ (système
accusatoire). Les États-Unis investissent massivement dans
$\mathcal{R}$ (recherche, normalisation technique). Singapour
optimise l'efficacité globale par ciblage plutôt que par
maximisation d'une dimension unique.

\cro{variable selon contexte}{variable selon priorité nationale}{variable selon système juridique}

\textbf{Conséquence pour le Cameroun~:} la question pertinente
n'est pas ``quel pays imiter'' mais ``quel arbitrage CRO correspond
aux priorités et contraintes camerounaises actuelles~?''
(Section~\ref{sec:recommandations-cameroun}).
\end{boxCRO}

% ============================================================================
\section{Recommandations pour le Cameroun}
\label{sec:recommandations-cameroun}
% ============================================================================

\subsection{Where Cameroun se situe aujourd'hui}

Le Chapitre~\ref{chap:droit-cam} a établi que le Cameroun dispose
d'un cadre légal récent et structuré (\loicam{2010/012},
\loicam{2024/017}) suivant la même logique de codification précise
que l'Allemagne. C'est un atout réel, souvent sous-estimé~: la
qualité du texte légal n'est pas le facteur limitant principal
identifié dans ce cours.

\begin{tableaucours}{p{4.0cm}p{7.5cm}}{%
    \tableauHeader{Atout actuel} &
    \tableauHeader{Élément}
}
Cadre légal codifié & \loicam{2010/012} et \loicam{2024/017}~:
    structure comparable aux textes allemands dans leur logique
    (codification précise plutôt que jurisprudentielle) \\
Coopération régionale & Accès aux mécanismes CEDEAO~/ AFRIPOL~/
    INTERPOL (Ch.~\ref{chap:norm-glob})~--- atout partagé avec
    les pays africains les plus avancés \\
Formation universitaire & Ce cours lui-même illustre l'émergence
    d'une formation structurée de niveau master, condition
    identifiée comme nécessaire dans tous les modèles étudiés \\
\end{tableaucours}
\vspace{-0.8em}
\begin{center}{\small\itshape Tableau~24.8 --- Atouts actuels du Cameroun}\end{center}

\subsection{Recommandations réalistes, ancrées dans les contraintes}

\begin{enumerate}
    \item \textbf{Cibler plutôt que disperser} (leçon de Singapour)~:
          identifier un ou deux domaines économiquement stratégiques
          (fraude Mobile Money, étant donné le poids économique de
          ce secteur au Cameroun) et y concentrer les premiers
          investissements de capacité spécialisée, plutôt que
          chercher une couverture exhaustive immédiate.
    \item \textbf{Maintenir la continuité institutionnelle}
          (leçon du Rwanda/Kenya)~: la progression observée dans
          les pays africains de référence résulte d'un effort
          soutenu sur~10--15~ans, pas de pics d'investissement
          ponctuels.
    \item \textbf{Investir dans la certification professionnelle}~:
          aucun des modèles étudiés n'a atteint une maturité
          forensique sans un système de qualification reconnu
          des experts (IACIS aux USA, UKAS au UK). Le
          Décret~69/DF/544 camerounais (Ch.~\ref{chap:droit-cam})
          est une base~; son application effective et son
          renforcement méritent une attention prioritaire.
    \item \textbf{Exploiter la coopération régionale comme levier},
          pas comme substitut~: la coopération AFRIPOL/CEDEAO
          (Ch.~\ref{chap:norm-glob}) compense des limites de
          capacité nationale, mais ne dispense pas de construire
          une expertise locale minimale et durable.
    \item \textbf{Rester lucide sur les comparaisons}~: le Cameroun
          ne dispose pas du budget de Singapour ni du volume
          d'affaires des États-Unis. Une stratégie réaliste se
          construit à partir des contraintes réelles, pas d'un
          objectif de ``rattrapage'' générique et non quantifié.
\end{enumerate}

\begin{boxalert}
\textbf{Une mise en garde contre les discours non étayés}

Les discours promettant un ``rattrapage rapide'' ou un
``leapfrogging technologique'' généralisé méritent une lecture
critique systématique~: ils s'appuient rarement sur une analyse
des contraintes réelles (budgétaires, humaines, institutionnelles)
et tendent à minimiser le temps nécessaire à la construction de
capacités durables. Les exemples étudiés dans ce chapitre
(Singapour~: 20--25~ans~; Rwanda/Kenya~: progression continue sur
10--15~ans) montrent que la maturité forensique se construit sur
le temps long, par des choix structurels cohérents et maintenus
--- pas par des annonces de rupture technologique.
\end{boxalert}

\separateur

\begin{boxresume}
\textbf{Ce qu'il faut retenir de ce chapitre~:}
\begin{itemize}
    \item \textbf{Aucun modèle national n'est universellement
          optimal}~: chaque système reflète un arbitrage CRO
          façonné par son contexte juridique propre (système
          accusatoire britannique~$\to$ priorité~$\mathcal{O}$~;
          investissement fédéral américain~$\to$ priorité
          $\mathcal{R}$ via la recherche).

    \item \textbf{Le Royaume-Uni} a développé les principes ACPO
          (Ch.~\ref{chap:meth-std}) sous la pression structurelle
          de l'obligation de \textit{disclosure}.

    \item \textbf{L'Allemagne} démontre qu'un système \textit{civil
          law} (comme le Cameroun) peut atteindre une maturité
          élevée par la codification précise, sans dépendre
          d'une jurisprudence accumulée~--- leçon directement
          transposable.

    \item \textbf{Singapour} illustre l'efficacité du ciblage
          sectoriel plutôt que la dispersion sur l'ensemble
          du spectre forensique~--- construit sur~20--25~ans.

    \item \textbf{Le Kenya et le Rwanda} montrent qu'un secteur
          économique moteur (fintech) ou une stratégie numérique
          nationale articulée peuvent accélérer significativement
          le développement de capacités spécialisées, sur une
          continuité de~10--15~ans.

    \item \textbf{Pour le Cameroun}~: le cadre légal
          (\loicam{2010/012}, \loicam{2024/017}) est un atout réel.
          Les recommandations réalistes privilégient le ciblage
          sectoriel (Mobile Money), la continuité institutionnelle,
          le renforcement de la certification professionnelle et
          l'exploitation de la coopération régionale comme levier.

    \item \textbf{Posture critique obligatoire}~: se méfier des
          discours de rattrapage technologique rapide non étayés
          par une analyse réaliste des contraintes et du temps
          nécessaire.
\end{itemize}
\end{boxresume}

\bigskip

\begin{boxexo}[title={\ding{45}\quad Exercice~24.1 --- Identifier les facteurs explicatifs \niveaubase}]
Pour chacun des pays suivants, identifiez (à partir du Tableau~24.7)
le facteur explicatif principal de sa force distinctive, et formulez
une question que vous poseriez pour vérifier si ce facteur est
transposable au Cameroun~:
\begin{enumerate}[(a)]
    \item Royaume-Uni
    \item Allemagne
    \item Singapour
    \item Kenya
\end{enumerate}
\end{boxexo}

\begin{boxexo}[title={\ding{45}\quad Exercice~24.2 --- Analyse critique d'un discours \niveaumoyen}]
Un consultant présente la proposition suivante~: ``Le Cameroun
doit investir massivement dans l'intelligence artificielle
forensique pour combler en~3~ans l'écart avec les pays occidentaux,
en s'appuyant sur l'absence de systèmes legacy qui freinent
l'innovation dans les pays développés.''
\begin{enumerate}[(a)]
    \item Identifiez dans cette proposition les éléments non
          étayés par les analyses de ce chapitre.
    \item Sur la base des exemples étudiés (Singapour, Rwanda,
          Kenya), quel délai réaliste suggéreriez-vous pour des
          progrès significatifs, et pourquoi~?
    \item Reformulez cette proposition de façon plus rigoureuse,
          en vous appuyant sur les recommandations de la
          Section~\ref{sec:recommandations-cameroun}.
\end{enumerate}
\end{boxexo}

\begin{boxexo}[title={\ding{45}\quad Exercice~24.3 --- Stratégie nationale argumentée \niveauexpert}]
Vous êtes consulté par l'ANTIC pour contribuer à l'élaboration
d'une feuille de route nationale de développement de la capacité
forensique numérique sur 10~ans.
\begin{enumerate}[(a)]
    \item Proposez un secteur prioritaire de ciblage initial
          (sur le modèle de Singapour) en justifiant votre choix
          par des éléments du contexte économique et criminel
          camerounais évoqués dans ce cours (Ch.~\ref{chap:gr-affaires},
          Ch.~\ref{chap:cas-int}).
    \item Identifiez trois éléments institutionnels (parmi ceux
          étudiés dans ce chapitre~: laboratoire centralisé,
          certification professionnelle, coopération régionale,
          formation universitaire) à prioriser dans les
          3~premières années, en justifiant l'ordre choisi.
    \item Appliquez la grille CRO à votre proposition~: quel
          arbitrage $\mathcal{C}$/$\mathcal{R}$/$\mathcal{O}$
          votre stratégie privilégie-t-elle, et est-ce cohérent
          avec les priorités identifiées en (a)~?
    \item Rédigez un paragraphe de mise en garde contre les
          écueils identifiés dans ce chapitre (dispersion des
          ressources, discours non étayés, rupture de continuité)
          à inclure dans votre feuille de route.
\end{enumerate}
\end{boxexo}

\bigskip

\begin{boxinfo}
\textbf{Transition.} Après cette mise en perspective comparative,
le cours se termine par un retour sur les fondements~: comment
synthétiser l'ensemble du parcours suivi depuis le Prologue, et
quelles sont les perspectives d'évolution de la discipline pour
les investigateurs camerounais formés aujourd'hui.
\end{boxinfo}